\section{Introduction}
\label{sec:intro}

On the Moon, an astronaut is a radio call away from flight controllers who watch every consumable and every step. On Mars that loop breaks. One-way light time to Earth ranges from about 3 to 22 minutes, and around each solar conjunction, every 26 months, when the Sun passes between Earth and Mars, it blocks contact for roughly two to three weeks \citep{mcbrayer2022comms}. NASA's exploration EVA concept of operations (ConOps) expects Mars crews to ``be more involved in planning their EVAs than as is currently done on ISS,'' using ``crew input, ground based planners, and automated planning tools'' under a delay of ``\textasciitilde4 and 20 minutes'' \citep{coan2020conops}. A recent review of planetary EVA tools finds that operational work has focused ``almost exclusively on microgravity rather than planetary surface operations'' \citep{elston2026review}.

Two gaps stand between today's tools and that need.
\emph{Gap~1: walks are routed from Earth.} EVA planning and monitoring are ground-based; NASA's EVA Mission System Software describes its role as ``empowering ground-based flight controllers'' \citep{miller2026emss}. From Mars, a question and its answer take up to 44 minutes. An onboard solar-particle alarm leaves only about 30 minutes to reach shelter \citep{lowe2025nowcast}. Earth therefore cannot re-route a walk in an emergency.
\emph{Gap~2: the planners on Mars are built for rovers.} The route planners that already run on Mars drive rovers \citep{carsten2009,toupet2026enav}. They model wheels, slip and slope, not a suited human's oxygen time, metabolic heat or walk back to shelter. Map portals show the data but plan nothing \citep{marstrek,mmgis}. Energy-optimal walking planners model effort, but not the full set of EVA constraints or guidance on Mars \citep{johnson2010sextant,penaasensio2024moonpath}.
Closing both gaps requires a crew-side tool that is
\textbf{R1}~\emph{Earth-independent}, planning and guiding on Mars without internet access;
\textbf{R2}~\emph{constraint-checked} against documented EVA limits, with ``do not go out'' always an acceptable answer;
\textbf{R3}~\emph{physiological}, following suited metabolism and suit heat rejection at Mars gravity;
\textbf{R4}~\emph{explainable} to a non-expert; and
\textbf{R5}~\emph{executable under link loss}, verified on arrival and still guiding the astronaut home without a link or GPS.

We present \Marswalk, built for the 2026 NASA International Space Apps Challenge (``Interplanetary Survival Guide: Martian Map''). In \Marswalk, the intravehicular (IV) crewmember in the habitat plans a walk on a 5\,km\,$\times$\,5\,km terrain model of Gale crater. The model fuses NASA and ESA orbital and rover data and includes Curiosity's 2025--26 boxwork campaign. The engine selects science stations, routes between them and checks the plan against twelve sourced constraints. A separate simulator then re-checks every plan before it is returned, and the engine explains the verdict. The confirmed plan is signed and sent as a package of about 11\,kB to the astronaut's phone-based augmented-reality (AR) navigator. The navigator estimates position without GPS, today from simulated drift and fixes, and keeps guiding when the link drops (\cref{fig:arch}). Our contributions are:
\begin{enumerate}[label=(\arabic*)]
  \item A formulation of crew-side Mars EVA planning as constrained orienteering on real terrain. Every constraint carries a source and an evidence class, and \stayin is a first-class result (\cref{sec:system}).
  \item A route engine that couples a suited metabolic model at Mars gravity with a suit heat-rejection governor. It computes one-to-many safety fields, selects stations with a labelled dynamic program that is exact over precomputed legs, re-checks plans in a separate time-stepped simulator, and explains verdicts counterfactually (\cref{sec:engine}).
  \item An Earth-independent execution path: a signed, compact route package with a conservative coarse safety grid, and a navigator with constant-time safety lookups that becomes autonomous on link loss. It is completed by a GPS-free positioning design and a suit hardware concept (\cref{sec:system,sec:interfaces,sec:hardware}).
  \item An evaluation built from public data. We compare the models with measured NASA data, quantify planning impact and parameter sensitivity over 460 simulated EVAs, cross-check verdicts against an independently written checker, and time the end-to-end system (\cref{sec:eval}).
\end{enumerate}
\Marswalk is a research prototype tested in simulation, not with suited people, real radios or EVA operators, and nothing in it is operational guidance. Its code, site-pack manifests and check scripts are open source.
